\chapter{CƠ SỞ LÝ THUYẾT}
\label{chap:lythuyet}

\section{SINH MÃ NGUỒN BẰNG MÔ HÌNH NGÔN NGỮ LỚN}
\lipsumtodo{LLM cho code, đánh giá bằng test, Pass@1; các bộ dữ liệu APPS \cite{hendrycks2021apps} và HumanEval \cite{chen2021humaneval}.
Quy ước thuật ngữ: dùng ``LLM'' làm thuật ngữ chung (khớp tên đề tài); ba mô hình của tiểu luận gọi là ``LLM mở cỡ nhỏ (7--9B tham số)''. Nêu một câu rằng một số tài liệu gọi nhóm mô hình cỡ này là small language model (SLM), nhưng ngưỡng SLM không thống nhất, nên tiểu luận không dựa vào ngưỡng nào và luôn ghi số tham số cụ thể.}

\section{BÀI TOÁN GIẢI BÀI VỚI BỘ KIỂM CHỨNG MỘT PHẦN}
\label{sec:baitoan}
\lipsumtodo{Phát biểu hình thức: yêu cầu $x$, lời giải $y$, bộ kiểm chứng một phần $V_{\text{pub}}$ (có phản hồi), bộ kiểm chứng ẩn $V_{\text{hid}}$, ngân sách $B$; mục tiêu tối đa xác suất $y$ vượt $V_{\text{hid}}$. Bảng ánh xạ hai trường hợp: sinh mã (public test / hidden test) và Text-to-SQL (chạy được, đúng lược đồ, đúng trên câu hỏi mẫu / đúng trên dữ liệu thật). Nêu hệ quả: khi nhiều lời giải cùng vượt $V_{\text{pub}}$, cần thêm tín hiệu để phân xử (critic).}

\section{TÌM KIẾM LÚC SUY LUẬN (TEST-TIME SEARCH)}
\subsection{Lấy mẫu nhiều lần và Best-of-N}
\lipsumtodo{}
\subsection{Monte Carlo Tree Search và công thức P-UCB}
\lipsumtodo{Selection, expansion, evaluation, backup (Max); P-UCB (Eq.~1--2 của bài báo).}

\section{RETHINKMCTS}
\subsection{Tìm kiếm ở mức thought}
\lipsumtodo{}
\subsection{Rethink tại lá và phản hồi mức block}
\lipsumtodo{CFG, trace biến, verbal feedback.}
\subsection{Dual reward có điều kiện}
\lipsumtodo{Critic chỉ chạy khi mọi public test đều pass (Eq.~3), thang $[-1,1]$.}

\section{KIỂM ĐỊNH THỐNG KÊ CHO SO SÁNH GHÉP CẶP}
\lipsumtodo{McNemar exact \cite{mcnemar1947}, hiệu chỉnh Holm \cite{holm1979}, khoảng tin cậy Wilson \cite{wilson1927}, bootstrap ghép cặp.}
